\documentclass[12pt,a4paper]{article}

\usepackage[utf8]{inputenc}
\usepackage[T1]{fontenc}
\usepackage[brazil]{babel}
\usepackage{geometry}
\usepackage{xcolor}
\usepackage{hyperref}
\usepackage{listings}

\geometry{a4paper,margin=2.5cm}
\hypersetup{colorlinks=true,linkcolor=blue,urlcolor=blue}

\definecolor{codegreen}{rgb}{0,0.45,0}
\definecolor{codegray}{rgb}{0.45,0.45,0.45}
\definecolor{codepurple}{rgb}{0.5,0,0.7}
\definecolor{backcolour}{rgb}{0.96,0.96,0.94}

\lstdefinestyle{mystyle}{
  backgroundcolor=\color{backcolour},
  commentstyle=\color{codegreen},
  keywordstyle=\color{magenta},
  numberstyle=\tiny\color{codegray},
  stringstyle=\color{codepurple},
  basicstyle=\ttfamily\footnotesize,
  breaklines=true,
  captionpos=b,
  keepspaces=true,
  numbers=left,
  numbersep=5pt,
  columns=fullflexible,
  showstringspaces=false,
  tabsize=2
}

\lstset{style=mystyle}

\title{\textbf{Vinheria IoT}\\Adega Climatizada da Vinheria Agnello}
\author{
  \textbf{Equipe CodeSquad-FIAP}\\
  Roger Viana Gonçalves de Alencar\\
  {\small RM 97540 --- Hardware e circuito}\\
  Kevin Benevides\\
  {\small Firmware Arduino e integração serial}
}
\date{Agosto de 2026}

\begin{document}

\maketitle
\newpage

\section{Objetivo e contextualização}

O projeto apresenta uma solução de Internet das Coisas para monitorar a adega
climatizada da Vinheria Agnello. A conservação de vinhos exige controle
ambiental constante: variações de temperatura aceleram o envelhecimento,
umidade inadequada prejudica as rolhas e a exposição luminosa pode provocar
degradação fotoquímica da bebida.

A solução coleta temperatura, umidade e luminosidade, transmite as medições em
JSON pela porta serial e permite que o Node-RED encaminhe os dados ao dashboard
e a um broker MQTT.

\section{Sensores e justificativas}

\subsection{Temperatura e umidade: DHT22}

O DHT22 foi selecionado por fornecer temperatura e umidade em um único sensor
digital. Para a conservação dos vinhos, a temperatura deve permanecer estável,
idealmente entre 12\(^{\circ}\)C e 18\(^{\circ}\)C, enquanto a umidade deve ficar
próxima de 70\%. Um ambiente seco pode retrair as rolhas e permitir a entrada
de oxigênio; um ambiente excessivamente úmido favorece fungos e danos aos
rótulos.

\subsection{Luminosidade: módulo LDR}

O módulo LDR mede a incidência de luz na adega. Exposição luminosa excessiva,
especialmente à radiação ultravioleta, pode desencadear reações nos compostos
de enxofre do vinho e produzir o defeito conhecido como \textit{light strike}.
O monitoramento permite detectar iluminação esquecida ou entrada indevida de
luz externa.

\section{Arquitetura de hardware}

O circuito é simulado no Wokwi com um Arduino Uno e possui a seguinte pinagem:

\begin{itemize}
  \item DHT22: alimentação em 5 V, terra comum e dados no pino digital D2;
  \item módulo LDR: alimentação em 5 V, terra comum e saída analógica no A0.
\end{itemize}

O arquivo executável do circuito está versionado em
\path{Arduino/VinheriaSensores/diagram.json}. A listagem completa aparece na
Seção~\ref{sec:fontes}.

\section{Firmware e formato dos dados}

O firmware inicializa o DHT22, configura o LDR e realiza uma amostragem a cada
4000 ms. O intervalo é controlado com \texttt{millis()}, sem bloquear o
microcontrolador com \texttt{delay()}.

Apesar de o circuito possuir dois sensores físicos, o DHT22 produz duas
medições. Assim, cada mensagem contém três campos semânticos:

\begin{lstlisting}[caption={Mensagem serial produzida pelo firmware}]
{"temperatura":14.2,"umidade":70.0,"luminosidade":512}
\end{lstlisting}

Temperatura é expressa em graus Celsius, umidade em percentual e luminosidade
como leitura bruta do conversor analógico do Arduino Uno, entre 0 e 1023. Se o
DHT22 falhar, temperatura e umidade são publicadas como \texttt{null}; desse
modo o documento continua sendo JSON válido e a leitura independente do LDR
não é perdida.

\section{Integração com Node-RED e MQTT}

O Arduino Uno utilizado não possui interface de rede. Portanto, a publicação
MQTT é responsabilidade do Node-RED, conforme o fluxo:

\begin{center}
Arduino Uno \(\rightarrow\) Serial \(\rightarrow\) Node-RED
\(\rightarrow\) HiveMQ \(\rightarrow\) Dashboard
\end{center}

A serial opera em 9600 baud, 8 bits de dados, sem paridade e um bit de parada.
Cada mensagem termina com \texttt{\textbackslash n}. No Node-RED, um nó serial
separa as linhas, um nó JSON converte a mensagem e os valores podem alimentar
gráficos ou ser publicados em um tópico MQTT.

Essa separação mantém o firmware responsável apenas pela aquisição e pelo
protocolo serial, enquanto conectividade, armazenamento e visualização ficam
na camada de integração.

\section{Reprodução e validação}

O diretório \path{Arduino/VinheriaSensores} reúne o sketch, o circuito e a
lista de bibliotecas. Para reproduzir a simulação:

\begin{enumerate}
  \item abrir um projeto Arduino Uno no Wokwi;
  \item copiar \path{VinheriaSensores.ino}, \path{diagram.json} e
        \path{libraries.txt};
  \item iniciar a simulação e abrir o monitor serial;
  \item alterar temperatura, umidade e luminosidade no simulador;
  \item confirmar uma linha JSON válida a cada quatro segundos.
\end{enumerate}

Para a compilação local, utiliza-se o Arduino CLI com o core AVR e a biblioteca
DHT da Adafruit.

\newpage
\section{Códigos-fonte}
\label{sec:fontes}

Os arquivos abaixo são incluídos diretamente no relatório para evitar que a
documentação fique diferente da versão executável.

\subsection{Esquema de hardware}

\lstinputlisting[caption={Circuito Wokwi em diagram.json}]{../../Arduino/VinheriaSensores/diagram.json}

\subsection{Firmware Arduino}

\lstinputlisting[language=C++,caption={Firmware VinheriaSensores.ino}]{../../Arduino/VinheriaSensores/VinheriaSensores.ino}

\end{document}
